\section{Polygonal Mesh Prossessing}

In this section, we 



\subsection{Triangulating Nonconvex Polyhedra}
\label{sec:triangulating}

A three-dimensional solid with polygonal faces and straight edges is called a \EMPH{polyhedra}. 
We often wish to decompose a polyhedra into tetrahedra much like
we decompose polygons into triangles.
Moreover, we would like to represent a polyhedra
with as few tetrahedra as possible \cite{simplify-mesh-1999}.

A polyhedron is \EMPH{simple} if it is homeomorphic to the sphere 
and the faces are all polygons.
Let $P$ be a a simple polyhedron with $n$  vertices.
In \cite{chazelle-lower-1984}, Chazell uses the Gauss-Bonnet
theorem to show that,
in the worst case, decomposing the $P$
 into tetrahedra requires $\Omega(n^2)$ tetrahedra.

 An edge $e$ in $P$ is
\EMPH{reflex} if the interior angle formed by its two incident faces
is greater than $\pi$. See \figref{reflex} for an example.
A vertex is reflex is it is incident to a reflex edge.
Let $r$ denote the number of reflex edges.


\begin{figure}[htb]
\centering
	\includegraphics[width=.3\textwidth]{meshes/reflex}
	\caption{A polyhedron with one reflex edge (blue), the two vertices incident to the blue
	edge are reflex vertices.}
\label{fig:reflex}
\end{figure}

Chazell and Palios give an
algorithm to triangulate a nonconvex simple polyhedra \cite{triangulating-polytope-1990}.
 Their algorithm creates a refinement with $O(n+r^2)$ tetrahedra
in $O(n\log r +r^2\log r)$ time and $O(n+r^2)$ space.
The algorithm first removes $n-4r$ non-reflex or flat vertices
to create a representation of the polyhedra with $O(r)$ vertices.
Then, vertical planes decompose the reduced polyhedra in to
$O(r^2)$ convex cells. They also establish that this bound is tight by exhibiting a class of polyhedra which cannot be decomposed into fewer than $n^2$ parts.

Chazelle and Shouraboura use the Gauss-Bonnet theorem to show any polyhedron
 of genus $g$ must have at least $g-1$ reflex edges~\cite{tetra-bounds-c-s-1994}.
 This implies that any polyhedron
can be decomposed with $\Theta(n+r^2)$ tetrahedra, regardless  of 
the genus!  We present their application.

In this application, to simplify computation, we scale the curvature of a vertex
by $\frac{1}{2\pi}$, so $K(v)=\frac{1}{2\pi}\left(2\pi-\sum_i \alpha_i\right)$,
We relate the number of reflex edges incident to a vertex $v$ to $K(v)$.

\begin{lemma}\label{lem:reflex-edge}
The number of reflex edges  incident to a vertex $v$  is at least $-K(v).$
\end{lemma}

\begin{proof}

For each vertex, center a sphere at the vertex,
see \figref{sphere-on-vert}.
The intersection of the polyhedron and the sphere
gives a `polygon' on the sphere with boundary $L$ consisting of great
circles. 
Scale the sphere to have unit radius. Then, the length of each arc
in $L$ is equal the the angle incident to $v$, so the total length of $L$ is
the sum of the angles incident to $v$, see \figref{sphere}.

Let $R$ be the number of reflex edges incident to $v$.
If $R$ is zero, then $P$ is convex at $v$ and has non-negative curvature
so $0>-K(v).$ If $R>0,$
reflex edges incident to $v$ correspond to reflex angles on $L$.
For each reflex angle, bisect it and reduce the 
number of reflex angles and obtain a decomposition
of the sphere into at most $R+1$ convex regions,
thus $L\leq 2\pi(R+1)$.

Thus, we have $L=\sum \alpha_f$ and the curvature 
$K(v)=1-\frac{1}{2\pi}\sum \alpha_f$ giving
$L=2\pi(1-K(v))$. Combining this with $L\leq 2\pi(R+1)$
gives $-K(v)\leq R.$


\end{proof}

\begin{figure}[htb]
        \centering
        \begin{subfigure}[b]{0.35\textwidth}
        \includegraphics[width=\textwidth]{meshes/reflex-vert-sphere}
        \caption{}
          \label{fig:sphere-on-vert}
        \end{subfigure}
          \hspace{1cm}
         \begin{subfigure}[b]{0.2\textwidth}
        \includegraphics[width=\textwidth]{meshes/sphere-intersection}
        \caption{}
        \label{fig:sphere}
        \end{subfigure}\\
		\caption{(a) A sphere centered at a reflex vertex. (b) The `polygon'
		of great circles where the polyhedron intersects the sphere. 
		\label{fig:sphere}}
\end{figure}


Next, we show
\begin{theorem}[Reflex Angles]\label{thm:reflex}

Any polyhedron of genus $g$ must have 
at least $g-1$ reflex dihedral angles. 

\end{theorem}
\begin{proof}
Let $r$ be the total number of reflex edges in a polyhedra.
By the classification of oriented surfaces, the Euler characteristic 
 is determined by the genus, $\chi=2-2g$.
By \lemref{reflex-edge}, $\sum_v -K(v)\leq 2 r$ since each
reflex edge is incident to two vertices.
By the Gauss-Bonnet theorem $\sum_vK(v)= 2-2g,$
so $-2r\leq 2-2g$ and $g\leq r+1$ as desired.

\end{proof}

Given a polyhedron of genus $g$, we can temporally 
duplicate the vertices around each essential cycle and insert disks,
 creating a polyhedron of genus zero, we then apply the algorithm to
decompose genus zero polyhedra given in \cite{triangulating-polytope-1990}.
Then remove the added disks.
Since we added $g$ disks the algorithm decomposes
the polyhedron into $O(n+ (r+g)^2)$ tetrahedra.
By \thmref{reflex}, $g\leq r+1$ showing that the upper bound
on the number of tetrahedra in a triangulation
of a polyhedra is $O(n+r^2)$ regardless of the genus.

\subsection{Digital Topology}
\label{sec:digital-topology}

In two dimensions, a black and white image consists of an array of pixels
where each pixel is either black or white.
In three dimensions,
instead of square pixels we have cubes called voxels.
Three dimensional images are use useful in medical images. 
For example, consider a
three dimensional images of the human heart \cite{bovik_handbook_2000}.
Given such an image, one might want to compute
the number of holes in the image to make sure it
matches the number of holes that are suppose to be in
the human heart. In this subsection, we share a result
of Chen and Rong where they 
use the Gauss-Bonnet theorem to 
compute genus of a digital image in $\R^3$ \cite{chen_digital_2010}.
In most of our other uses of the Gauss-Bonnet theorem we begin with surface where we know
the Euler characteristic, the sphere or the disk, and use the theorem to tell us
about the local curvature. In this application, we are going in the opposite direction,
we are using the local curvature to compute the genus.


Digital images are commonly reffered to as \emph{cubical complexes},
where each cube has side length one \cite{kong_digital_1989}.
An \emph{elementary cube} $Q$ is a product
$$Q=I_1\times I_2 \times I_3$$
where each $I_i$ is a unit interval of the form $I=[\ell,\ell+1]$
called a \emph{coordinate}.
We will often associate each elementary cube with a lattice point
see \figref{4-square-and-dual} for an example.
%In many cases, given a cubical complex, one is interested in having
%a computer to to classify the image.
%Operations toward image classification include
%object counting, border following and computing 
%the number of holes \cite{kong_digital_1989}.

Our goal is to use the Gauss-Bonnet theorem to compute
the genus of a three dimensional cubical complex.
The definition of genus, \defref{genus}, is given in terms
of simple closed curves. In order to describe a closed curve in a cubical
complex we first need to resolve an interesting paradox.
In two dimensions, it is not clear what it means
for two pixels to be adjacent. For example, in \figref{4-square}
we see four black pixels. Should these pixels represent a closed
curve? If we define the pixels adjacent to a given pixel to be the 
pixels to the left, right, above, and below, then each black pixel isolated
and the four pixels should not form a closed curve. However,
with adjacency defined this way, the white pixel in the middle is completely surrounded 
and the four black pixels should represent a closed curve because
they separate the white pixels into an inside and an outside.
In order to resolve this paradox introduce different notions of adjacency.


In two dimensions, two points are \emph{eight-adjacent}
if they are distinct and each coordinate of one differs from
the other by at most one.
Two points are \emph{four-adjacent} if
they are eight-adjacent and differ in at most one of their coordinates.
Consider four vertices adjacent to a single vertex $v$.

Returning to our paradox in \figref{4-square}.
If we use eight-adjacency
for the white vertices and four-adjacency for the black.
Then the four black pixels do not form a closed curve 
and the center white pixel is not enclosed, so we avoid our paradox. 
Notice we could have chosen four adjacency
for the white vertices and eight-adjacency for the black pixels and 
now the four black pixels do form a closed curve and the white pixel
is surrounded and again we do not have a paradox.



We make a similar choice in three dimensions, two points
are \emph{26-adjacent} if they are distinct and each coordinate 
entry differs by at most one.
The points are \emph{18-adjacent} if they are 26-adjacent
and different in  a most two of their coordinates.
The points are \emph{6-adjacent} if they are 18-adjacent and differ 
in at most one of their coordinates.
If a set of points $S$ lattice points cannot be
partitioned into two subsets that are not
$n$-adjacent is \EMPH{$n$-connected}.



\begin{figure}[htb]
        \centering
        \begin{subfigure}[b]{0.25\textwidth}
        \includegraphics[width=\textwidth]{digital/4-square}
        \caption{}
          \label{fig:4-square}
        \end{subfigure}
          \hspace{.3cm}
         \begin{subfigure}[b]{0.25\textwidth}
        \includegraphics[width=\textwidth]{digital/4-square-dual}
        \caption{}
        \label{fig:4-square-dual}
        \end{subfigure}\\
		\caption{(\subref{fig:4-square}) A two dimensional image with four black pixels.
		Do the black pixels represent a closed curve? (\subref{fig:4-square-dual}) The same image
		represented on the lattice.
		\label{fig:4-square-and-dual}}
\end{figure}

A \EMPH{digital picture} is a quadruple $(V,m,n,B)$ where
$V=\R^2$ and $(m,n)=(4,8)$ or $V=\R^3$ and $(m,n)=(6,26)$
and $B$ is a subset of $V$. Elements of $B$ are black vertices
and elements not in $V$ are white.

\begin{figure}[htb]
        \centering
        \includegraphics[width=.25\textwidth]{digital/26-neighbors}
		\caption{A vertex in a three dimensional cubical complex and its 26 neighbors.
		\label{fig:26-neighbors}}
\end{figure}

One may ask, why don't we triangulate the cubical complex and
count vertices, edges, and faces to determine the genus.
One reason against doing this is that 
decomposing a cubical complex into a simplicial complex
results in 24 times as many highest dimensional cells \cite{Kaczynski2003}.


The objects we wish to compute the genus of are closed, orientable, cubical complexes in $\R^3$ with $(6,26)$-connectivity.
If $M$ is such a complex, there are
six types of digital surface points, these are shown in \figref{6-types}.


\begin{figure}[htb]
        \centering
        \begin{subfigure}[b]{0.2\textwidth}
        \includegraphics[width=\textwidth]{digital/3-neighbors}
        \caption{Three neighbors.}
          \label{fig:3-neighbors}
        \end{subfigure}
          \hspace{.3cm}
         \begin{subfigure}[b]{0.25\textwidth}
        \includegraphics[width=\textwidth]{digital/4-neighbors-flat}
        \caption{Four neighbors.}
        \label{fig:4-neighbors-flat}
        \end{subfigure}
         \hspace{.3cm}
         \begin{subfigure}[b]{0.2\textwidth}
        \includegraphics[width=\textwidth]{digital/4-neighbors-bent}
        \caption{Four neighbors too.}
        \label{fig:4-neighbors-bent}
        \end{subfigure}\\
         \begin{subfigure}[b]{0.25\textwidth}
        \includegraphics[width=\textwidth]{digital/5-neighbors-a}
        \caption{Five neighbors.}
          \label{fig:1-neighbors}
        \end{subfigure}
          \hspace{.2cm}
         \begin{subfigure}[b]{0.25\textwidth}
        \includegraphics[width=\textwidth]{digital/5-neighbors-b}
        \caption{Six neighbors.}
        \label{fig:4-neighbors-flat}
        \end{subfigure}
         \hspace{.2cm}
         \begin{subfigure}[b]{0.25\textwidth}
        \includegraphics[width=\textwidth]{digital/6-neighbors}
        \caption{Six neighbors too.}
        \label{fig:6-neighbors}
        \end{subfigure}
		\caption{This six types of digital surface points.
		\label{fig:6-types}}
\end{figure}

Let $M_i$ denote the set of digital points with $i$ neighbors and $K_i$
the curvature.
Then, by \defref{discrete-curvature-vertex}, we compute the discrete curvature
at each of this points to obtain
\begin{enumerate}[(a)]
\item $K_3=\pi/2,$
\item $K_4=0,$
\item $K_5=-\pi/2,$
\item $K_6=-\pi.$
\end{enumerate}

Then, by the Gauss-Bonnet theorem, we have
$$\sum_{i=3}^6K_i |M_i|=2\pi (2-2g).$$
A linear time algorithm to compute the genus is the following.
Iterate through all points in $M$ and count the neighbors at each point
and keep track of $M_i$. Then calculate the genus
by rearranging the above equation to obtain
$$g=1+(|M_5|+2|M_6|-|M_3|)/8.$$

\subsection{Homotopy Testing}
\label{sec:homotopy}

Homotopy classes of curves are important  because...
In this section, we share applications
where the Gauss-Bonnet theorem is used to determine
if curve on a surface is contractible.
An explanation of this algorithm is given in a video lecture 
by Jeff Erickson \cite{erickson-lecture}.

\begin{figure}[htb]
\centering
\includegraphics[width=.3\textwidth]{homotopy/contractable}
\caption{The blue curve is not contractable, the red curve is contractable.}
\label{fig:contractable}
\end{figure}

In this application, we will use a modified version
of the Gauss-Bonnet theorem. Given a combinatorial surface
$S$ define the curvature of a face to be $k(f)=1-\sum \beta_f$,
$k(v)=1-\frac{1}{2}\deg(v)+\sum_{v\in V} \alpha_v$.
Summing over all faces and vertices in $S$ gives,
\begin{theorem}[Modified Gauss-Bonnet]\label{thm:modify-g-b}
For a combinatorial surface $S$,
$\sum_{f\in F} k(f)+\sum_{v\in V}k(v)=\chi(S).$
\end{theorem}

\begin{proof}

$$\sum_{f\in F} k(f)+\sum_{v\in V}k(v)=|F|-\sum_{v\in V}\alpha_v + |V|-|E|+\sum_{v\in V}\alpha_v=\chi(S).$$
\end{proof}


Here, a we have a combinatorial surface that we call a map denoted $S$
with genus $g\geq 2.$ The case $g=0$ is the sphere where every curve
is contractible and when $g=1$ contractibility can be determined by a counting
argument\todo{explain}.
A curve is a closed walk, given as alternating sequence of vertices
and edges in $S$.
A \EMPH{homotopy} between two closed curves $\gamma_1$ and $\gamma_2$ that 
share a point $p_0$ is a continuous map $H\colon [0,1]\times \Sp^1 \to \mathbb{R}^2$ 
such that $H(0,\cdot)=\gamma_1$, $H(1,\cdot)=\gamma_2$, and $H(s,0)=p_0=H(s,1)$.

On a combinatorial surface, homotopies can be decomposed
into discrete moves called edge spikes, edge unspikes  and face flips.
The question we consider is: Given a closed walk $W$ in a map $\Sigma$ is there a finite
sequence of moves that reduces the curve to a trivial walk?


First, we transform $S$ into a simpler object called a system of loops.  
Let $T$ be a spanning tree of $S$, let $C$ be the edges in a co-tree
(a spanning tree of the dual), and let $L$ denote the left over edges,
Each edge in $S$ is in one of $(T,L,C)$.
Contract edges in $T$, delete edges in $C$ a system of loops $\Delta$.
When we contract and delete edges in $S$ we might affect edges in our walk $W$.
If we contract an edge in $W$, delete the edge from $W$.
If we delete and edge $e\in W$, face flip to avoid the deleted edge.
Get a new walk $W' \in \Delta$ homotopic to $W\in \Sigma$
now only haveing one vertex. Then, $W'$ has one vertex
with degree $4g$, one face with $4g$ edges, and $2g$ loops.
The length of the walk increases by at most $2g$.


The universal covering space of $S$ denoted $U$ is a plane
and when $g\geq 2$ $U$ has a natural hyperbolic geometry,
tiled by $4g$-gons.


Big idea: a curve is contractible closed walk in $S$
if and only if the walk is closed $U$.
Thus, our problem is equivalent to the following: Given a walk
in the universal cover of a system of loops, is it closed?

Here the walk is given as a starting vertex then we have a list of which
edges to take at each intermediate vertex. 
All vertices look the same, so we must determine if we have ended where
we began. 

Look for spurs or taking and instances where we talk the long way around a face, 
shorten the walk.
Any nontrivial contractible cycle contains either a spur or a bracket \cite{gertsen-short-1990}.
\begin{lemma}[Dehn's Lemma]\label{lem:dehn}
If $g\geq 2,$ then any nontirival closed walk has either a spur
or $4g-2$ consecutive edges on the boundary of a face.
\end{lemma}
\begin{proof}
Sketch: 
Any nontrival closed walk bounds a disk $D$ with $\chi(D)=1$, by the Gauss-Bonnet theorem
$$\sum_v k(v)+\sum_f k(f) =1.$$

Each face has $4g$ edges, each internal vertex has degree $4g$
and each  boundary vertex  has degree less than $4g$.
The angle of each angle on a face  is $\frac{1}{4}$,
thus, $k(f)=1-g<0,$ for internal vertices  $k(v)=1-g<0$ and for all
vertices on the boundary, $k(v)=\frac{3}{4}-\frac{\deg(v)}{4}$.
Boundary vertices fall into three  categories: convex, where $k(v)=\frac{1}{4}$
flats,  where $k(v)=0$  and concave where $k(v)<0$.

By G-B, $$|F|(1-g)+|v_{convex}|\frac{1}{4}\geq 1$$
and 
$$|v_{convex}|\geq (4g-4)|F|+4.$$

We divide by $|F|$ to determine the average number
of convex vertices per face to be greater than $4g-4$.
Thus, there exists some face that has $4g-3$ consecutive edges in the walk.
\end{proof}

This gives an algorithm for determining if a walk is closed in the universal
covering space.
Look for spurs and long boundary subpaths, the walk is closed if and  only if
we can  shrink the curve.
Label edges, walk is a sequence of labels
look at intervals of $4g-2$ in a walk, $8g$ paths
that represent long boundary paths, $4g$ spurs.
Slide window and look for spurs or long boundary paths.
If you find one remove it.

Brute force $O(g^3\ell)$ overall.
Can speed it up with Erickson DFA  idea
$O(g^2+g\ell)$.
Overall runtime $O(n+g^2+g\ell)$ time.

In trouble if $g$is big. Erickson uses system of quads,
radial map, $O(n)$ runtime  \cite{erickson-whittlesey-2013}.








\noindent \textbf{Exercises}


\begin{enumerate}
	\item 
	
	\item 
	
	\item 
	
\end{enumerate}

\pagebreak
